\documentclass[12pt,a4paper]{article}
\usepackage[utf8]{inputenc}
\usepackage[T1]{fontenc}
\usepackage{geometry}
\usepackage{enumitem}
\usepackage{fancyhdr}
\usepackage{lastpage}
\usepackage{hyperref}

\geometry{margin=2.5cm}
\pagestyle{fancy}
\fancyhf{}
\rhead{Page \thepage\ of \pageref{LastPage}}
\lhead{LNP201 - Quiz (en)}
\renewcommand{\headrulewidth}{0.4pt}

\title{Quiz: LNP201}
\author{Language: en}
\date{\today}

\begin{document}

\maketitle
\thispagestyle{fancy}

\section*{Instructions}
Answer all questions by selecting the best option (A, B, C, or D). The answer key is provided at the end of this document.

\bigskip


\subsection*{Question 1}
What is the main goal of the Lightning Network?

\begin{enumerate}[label=\Alph*., leftmargin=*]
\item Eliminate transaction fees on the Bitcoin blockchain
\item To improve Bitcoin's scalability to enable fast and low-cost payments
\item Create an alternative to Bitcoin for high-value payments
\item Simplify exchanges between different cryptocurrencies
\end{enumerate}

\bigskip


\subsection*{Question 2}
In a payment channel, what is the capacity of the channel?

\begin{enumerate}[label=\Alph*., leftmargin=*]
\item The total amount of funds that the two participants can exchange
\item The time limit for which the channel can remain open
\item The maximum amount one participant can send in a single transaction
\item The number of transactions a channel can support per day
\end{enumerate}

\bigskip


\subsection*{Question 3}
If Alice sends 40,000 satoshis to Bob from a payment channel where she initially has 100,000 satoshis, how much funds does she have left after the transaction?

\begin{enumerate}[label=\Alph*., leftmargin=*]
\item 40,000 satoshis
\item 30,000 satoshis
\item 140,000 satoshis
\item 60,000 satoshis
\end{enumerate}

\bigskip


\subsection*{Question 4}
What is the main rule regarding the funds a participant can send in a payment channel?

\begin{enumerate}[label=\Alph*., leftmargin=*]
\item They can send up to half of the total capacity of the channel.
\item They can send any amount, as long as there is a prior agreement.
\item They can borrow funds from the other participant's side if needed.
\item They can only send the funds from their own side of the channel.
\end{enumerate}

\bigskip


\subsection*{Question 5}
What does the Lightning Network allow, compared to traditional transactions on the Bitcoin blockchain?

\begin{enumerate}[label=\Alph*., leftmargin=*]
\item Transactions between different blockchains
\item Anonymous and untraceable transactions
\item Transactions in large quantities without volume limits
\item Fast transactions with minimal fees
\end{enumerate}

\bigskip


\subsection*{Question 6}
What is the primary role of a private key in the Bitcoin system?

\begin{enumerate}[label=\Alph*., leftmargin=*]
\item Serve as an identifier to receive bitcoins
\item Create locking scripts to protect bitcoins
\item Sign transactions to satisfy a script
\item Store funds securely on the blockchain
\end{enumerate}

\bigskip


\subsection*{Question 7}
What is a UTXO in the context of Bitcoin?

\begin{enumerate}[label=\Alph*., leftmargin=*]
\item A piece of unspent bitcoins associated with a script
\item An amount of bitcoin divided into multiple parts
\item A fully confirmed transaction on the blockchain
\item A public key that generates Bitcoin addresses
\end{enumerate}

\bigskip


\subsection*{Question 8}
Why are UTXOs often compared to banknotes?

\begin{enumerate}[label=\Alph*., leftmargin=*]
\item They represent small fixed amounts of bitcoins
\item They can be used multiple times until depleted
\item They must be spent in their entirety
\item They are only usable in certain specific transactions
\end{enumerate}

\bigskip


\subsection*{Question 9}
How does a 2/2 multisignature address work in Bitcoin?

\begin{enumerate}[label=\Alph*., leftmargin=*]
\item It requires two signatures to unlock and spend the funds.
\item It is used only to send funds to public addresses.
\item It authorizes transactions only in specific cases of locking.
\item It allows spending the funds with just one of the two keys.
\end{enumerate}

\bigskip


\subsection*{Question 10}
In a Bitcoin transaction, what does the locking script do?

\begin{enumerate}[label=\Alph*., leftmargin=*]
\item It imposes specific conditions to spend the associated funds
\item It automatically sends the funds to multiple addresses
\item It temporarily locks the funds to prevent excessive spending
\item It determines the value of the UTXO to be created
\end{enumerate}

\bigskip


\subsection*{Question 11}
What is the first step in opening a Lightning channel between two participants?

\begin{enumerate}[label=\Alph*., leftmargin=*]
\item Sign the withdrawal transaction before the exchange of public keys
\item Create a 2/2 multisignature address without prior exchange
\item Immediately publish a deposit transaction on the blockchain
\item Exchange messages containing the amounts and public keys
\end{enumerate}

\bigskip


\subsection*{Question 12}
Why does Alice create a withdrawal transaction before publishing the deposit transaction?

\begin{enumerate}[label=\Alph*., leftmargin=*]
\item To allow Bob to spend the funds directly
\item To increase the channel's capacity with each transaction
\item To perform a Lightning transaction immediately after opening
\item To retrieve her funds in case of a problem with Bob
\end{enumerate}

\bigskip


\subsection*{Question 13}
When is a Lightning channel officially open?

\begin{enumerate}[label=\Alph*., leftmargin=*]
\item When the deposit transaction is confirmed on the Bitcoin blockchain
\item As soon as the public keys are exchanged between participants
\item When the withdrawal transaction is signed by both parties
\item As soon as the 2/2 multisignature address is created
\end{enumerate}

\bigskip


\subsection*{Question 14}
What happens if Alice does not obtain Bob's signature for the withdrawal transaction?

\begin{enumerate}[label=\Alph*., leftmargin=*]
\item She can no longer communicate with Bob on the Lightning Network
\item She risks losing access to her funds deposited in the multisignature address
\item The channel opens with a limited capacity
\item Bob automatically receives the funds in the channel
\end{enumerate}

\bigskip


\subsection*{Question 15}
What is the last step to open a channel on the Lightning Network?

\begin{enumerate}[label=\Alph*., leftmargin=*]
\item Send a confirmation message of the opening to the other party
\item Verify the balance of funds of each participant
\item Sign a withdrawal transaction to recover the funds
\item Publish the deposit transaction on the Bitcoin blockchain
\end{enumerate}

\bigskip


\subsection*{Question 16}
What does a commitment transaction in a Lightning channel reflect?

\begin{enumerate}[label=\Alph*., leftmargin=*]
\item The total number of transactions made
\item The amount of fees to be paid for the transaction
\item The new distribution of funds between the two parties
\item The opening of a new payment channel
\end{enumerate}

\bigskip


\subsection*{Question 17}
When is a commitment transaction published on the Bitcoin blockchain?

\begin{enumerate}[label=\Alph*., leftmargin=*]
\item Eventually, when the channel is closed
\item Immediately after each transfer in the channel
\item Before each new commitment transaction
\item When the available funds are exhausted
\end{enumerate}

\bigskip


\subsection*{Question 18}
What happens if Bob sends an invoice to Alice for 30,000 satoshis?

\begin{enumerate}[label=\Alph*., leftmargin=*]
\item The invoice is paid through an on-chain transaction on Bitcoin.
\item Alice creates a commitment transaction to transfer 30,000 satoshis to Bob in the channel.
\item Bob automatically receives the 30,000 satoshis in his wallet.
\item Alice must open a new channel to make this transfer.
\end{enumerate}

\bigskip


\subsection*{Question 19}
How is a new channel state created after each transfer?

\begin{enumerate}[label=\Alph*., leftmargin=*]
\item By posting an on-chain transaction after each transfer
\item By increasing the total capacity of the channel
\item By closing and reopening the channel after each transaction
\item By creating a new commitment transaction with the updated distribution
\end{enumerate}

\bigskip


\subsection*{Question 20}
What does the commitment transaction system guarantee to participants of a Lightning channel?

\begin{enumerate}[label=\Alph*., leftmargin=*]
\item That the fees are shared between the two participants
\item That the channel's capacity increases with each transaction
\item That transactions are visible on the Bitcoin blockchain
\item That each can recover their funds by publishing the last signed transaction
\end{enumerate}

\bigskip


\subsection*{Question 21}
What is the main risk in commitment transactions on the Lightning Network?

\begin{enumerate}[label=\Alph*., leftmargin=*]
\item Forgetting to sign the commitment transaction
\item Accidentally sending funds to another channel
\item Publishing an old transaction to reclaim more funds
\item Blocking the transaction with an inappropriate timelock
\end{enumerate}

\bigskip


\subsection*{Question 22}
How does the timelock protect against cheating in commitment transactions?

\begin{enumerate}[label=\Alph*., leftmargin=*]
\item It cancels the transaction in case of an attempted cheat
\item It blocks the use of the revocation key
\item It prevents the offending party from immediately retrieving their funds
\item It automatically refunds the funds to the cheater
\end{enumerate}

\bigskip


\subsection*{Question 23}
What does the revocation key allow to do in case of cheating?

\begin{enumerate}[label=\Alph*., leftmargin=*]
\item Cancel the commitment transaction
\item Double the channel's capacity
\item Recover all the funds by punishing the cheater
\item Prevent any further transactions in the channel
\end{enumerate}

\bigskip


\subsection*{Question 24}
Why do Alice and Bob exchange their secrets for previous commitment transactions?

\begin{enumerate}[label=\Alph*., leftmargin=*]
\item To allow each to protect themselves if the other publishes an old transaction
\item To divide the funds equally
\item To increase the channel's capacity
\item To block each other's access to the funds
\end{enumerate}

\bigskip


\subsection*{Question 25}
What is the combined effect of the timelock and the revocation key on commitment transactions?

\begin{enumerate}[label=\Alph*., leftmargin=*]
\item They force the opening of a new channel if the transaction is incorrect
\item They reduce the channel's capacity in case of cheating
\item They deter participants from cheating by offering an immediate punishment
\item They automatically lock the funds of the guilty party
\end{enumerate}

\bigskip


\subsection*{Question 26}
What are the three methods of closing a channel on the Lightning Network?

\begin{enumerate}[label=\Alph*., leftmargin=*]
\item Direct closure, indirect closure, closure with claim
\item Rapid closure, slow closure, closure with dual signature
\item Cooperative closure, forced closure, closure with cheating
\item Secure closure, automated closure, closure with loss of funds
\end{enumerate}

\bigskip


\subsection*{Question 27}
Which channel closing method is the most economical and fast?

\begin{enumerate}[label=\Alph*., leftmargin=*]
\item Cooperative closure
\item None are more economical and fast
\item Closure with cheating
\item Forced closure
\end{enumerate}

\bigskip


\subsection*{Question 28}
Why is forced closure less desirable?

\begin{enumerate}[label=\Alph*., leftmargin=*]
\item It forces the recreation of a channel later
\item It requires authorization from Satoshi Nakamoto
\item It is slower, and the fees may be inappropriate
\item It prevents the use of the revocation key
\end{enumerate}

\bigskip


\subsection*{Question 29}
What can Bob do if he detects an attempt by Alice to cheat during the channel closure?

\begin{enumerate}[label=\Alph*., leftmargin=*]
\item Permanently close the channel with no option to reopen
\item Use the revocation key to recover all the funds
\item Cancel Alice's commitment transaction
\item Claim only half of the available funds
\end{enumerate}

\bigskip


\subsection*{Question 30}
In what situation is a forced closure of a channel used?

\begin{enumerate}[label=\Alph*., leftmargin=*]
\item When the time lock is insufficient
\item When the capacity of the channel needs to be increased
\item When transaction fees are too high
\item When one of the parties no longer responds
\end{enumerate}

\bigskip


\subsection*{Question 31}
How does the Lightning Network allow sending funds to a node that is not directly connected?

\begin{enumerate}[label=\Alph*., leftmargin=*]
\item By sending the funds directly with multisignature addresses
\item By routing the payment through intermediary nodes
\item By using on-chain transactions between each node
\item By automatically establishing a temporary channel
\end{enumerate}

\bigskip


\subsection*{Question 32}
What limits the amount of an indirect payment on the Lightning Network?

\begin{enumerate}[label=\Alph*., leftmargin=*]
\item The total capacity of the Lightning Network
\item The routing fees set by each node
\item The number of intermediary nodes involved
\item The smallest liquidity available in a channel of the chosen path
\end{enumerate}

\bigskip


\subsection*{Question 33}
What is the role of routing fees in the Lightning Network?

\begin{enumerate}[label=\Alph*., leftmargin=*]
\item To limit the number of transactions on the network
\item To incentivize intermediary nodes to provide their liquidity for the transfer of payments
\item To ensure that each channel has sufficient funds
\item To automatically calculate the shortest path
\end{enumerate}

\bigskip


\subsection*{Question 34}
What does onion routing in the Lightning Network entail?

\begin{enumerate}[label=\Alph*., leftmargin=*]
\item To send funds to all nodes on the path simultaneously
\item To calculate fees for each intermediate node independently
\item It means that each intermediate node only knows its immediate neighbors on the path.
\item To directly transmit funds from node to node knowing the final recipient
\end{enumerate}

\bigskip


\subsection*{Question 35}
What information must the sending node know to perform an onion routing payment?

\begin{enumerate}[label=\Alph*., leftmargin=*]
\item The capacity of each channel
\item The cryptographic signature of each node on the path
\item The network topology and the complete route to the recipient
\item Only the next intermediate node
\end{enumerate}

\bigskip


\subsection*{Question 36}
What is the main goal of HTLCs in the Lightning Network?

\begin{enumerate}[label=\Alph*., leftmargin=*]
\item To reduce transaction fees for each payment through intermediary nodes
\item To establish direct connections between all network nodes with their trust
\item To enable secure payment transfers via intermediary nodes without requiring trust in them
\item To increase the capacity of each payment channel
\end{enumerate}

\bigskip


\subsection*{Question 37}
What condition must be met for the recipient of an HTLC to receive the funds?

\begin{enumerate}[label=\Alph*., leftmargin=*]
\item They must provide a secret (preimage) to unlock the payment.
\item They must accept the routing fees of each intermediary node.
\item They must wait for the timelock to expire before retrieving the funds.
\item They must sign a new commitment transaction.
\end{enumerate}

\bigskip


\subsection*{Question 38}
What happens if a payment through HTLC is not completed within the allotted time?

\begin{enumerate}[label=\Alph*., leftmargin=*]
\item The payment is canceled, and the funds are returned to the sender.
\item The payment is kept by the last intermediary node
\item The funds are distributed to all nodes in the path
\item The funds are permanently blocked
\end{enumerate}

\bigskip


\subsection*{Question 39}
What role does the timelock play in the operation of HTLCs?

\begin{enumerate}[label=\Alph*., leftmargin=*]
\item It provides a timeframe to complete the payment before cancellation
\item It adds a layer of privacy to the transaction
\item It applies additional fees for each intermediary node
\item It increases the capacity of the channel used
\end{enumerate}

\bigskip


\subsection*{Question 40}
What happens if a channel with an active HTLC is closed forcibly?

\begin{enumerate}[label=\Alph*., leftmargin=*]
\item The funds are redistributed to the previous nodes.
\item The HTLC is automatically cancelled.
\item The HTLC is included in the commitment transactions as an output.
\item The channel cannot be closed until the HTLC has expired.
\end{enumerate}

\bigskip


\subsection*{Question 41}
What is the main role of the sending node in routing on the Lightning Network?

\begin{enumerate}[label=\Alph*., leftmargin=*]
\item Communicate directly with all intermediary nodes
\item Store the liquidity information of all the network's channels
\item Determine the transaction fees for each intermediary node
\item Calculate the complete route to the recipient
\end{enumerate}

\bigskip


\subsection*{Question 42}
Why is it difficult for a transmitting node to estimate the exact liquidity of a channel?

\begin{enumerate}[label=\Alph*., leftmargin=*]
\item Because the total capacity of the channels is constantly changing
\item Because the fees of intermediary nodes are not published
\item Because channels automatically close after each transaction
\item Because only the two participants know the precise distribution of funds
\end{enumerate}

\bigskip


\subsection*{Question 43}
What criterion can help a node choose an efficient payment route?

\begin{enumerate}[label=\Alph*., leftmargin=*]
\item The geographical distance between the nodes
\item The probability of success based on the total capacity of the channels
\item The number of payments made by the target node
\item The presence of nodes with a low fee level
\end{enumerate}

\bigskip


\subsection*{Question 44}
What information can Bob include in an invoice to facilitate routing?

\begin{enumerate}[label=\Alph*., leftmargin=*]
\item The complete map of all the network's channels
\item The private key of intermediate nodes
\item Nearby channels with sufficient liquidity or the existence of private channels
\item The exact amount of available liquidity for each channel
\end{enumerate}

\bigskip


\subsection*{Question 45}
What does a sending node do when a payment route fails due to insufficient liquidity?

\begin{enumerate}[label=\Alph*., leftmargin=*]
\item It automatically sends a refund request to the recipient
\item It tries another route until it finds a functional one
\item It closes the channel and creates a new direct channel
\item It asks the intermediary nodes to increase their liquidity
\end{enumerate}

\bigskip


\subsection*{Question 46}
What does the LNURL-withdraw feature allow users to do?

\begin{enumerate}[label=\Alph*., leftmargin=*]
\item Make a direct payment without routing fees
\item Connect their Lightning wallet without a password
\item Withdraw funds from a service by scanning a QR code
\item Send funds without recipient confirmation
\end{enumerate}

\bigskip


\subsection*{Question 47}
In a Lightning invoice, what does the prefix lnbc in the human-readable part mean?

\begin{enumerate}[label=\Alph*., leftmargin=*]
\item That the invoice is valid only on the Bitcoin testnet
\item It indicates a Lightning transaction on the Bitcoin blockchain
\item That it is an on-chain transaction
\item That the invoice is encoded in bech32
\end{enumerate}

\bigskip


\subsection*{Question 48}
What is the role of LNURL in Lightning transactions?

\begin{enumerate}[label=\Alph*., leftmargin=*]
\item Facilitate the creation of channels for each payment
\item To simplify the interaction between nodes and Lightning applications
\item Ensure the confidentiality of payment information
\item Automatically calculate transaction fees
\end{enumerate}

\bigskip


\subsection*{Question 49}
How does the Keysend technology enable a fund transfer without an invoice?

\begin{enumerate}[label=\Alph*., leftmargin=*]
\item The funds are directly credited without additional security
\item The sender generates a secret in the HTLCs to unlock the funds
\item The recipient generates a hash to validate the transfer
\item The payment is validated through a direct channel between sender and recipient
\end{enumerate}

\bigskip


\subsection*{Question 50}
What is the primary function of a Lightning invoice?

\begin{enumerate}[label=\Alph*., leftmargin=*]
\item Send funds without an on-chain transaction
\item Activate HTLCs on a channel
\item Request a payment from the sending node
\item Open a new Lightning channel
\end{enumerate}

\bigskip


\subsection*{Question 51}
Why might a seller make a fictitious payment in a channel?

\begin{enumerate}[label=\Alph*., leftmargin=*]
\item To decrease routing fees
\item To move liquidity to the opposite side
\item To increase the total capacity of their channels
\item To lock funds in the long term
\end{enumerate}

\bigskip


\subsection*{Question 52}
What does the triangle opening technique allow participating nodes to achieve?

\begin{enumerate}[label=\Alph*., leftmargin=*]
\item A direct link with every other network node
\item An unlimited number of transactions without confirmation
\item Reduced routing fees
\item Incoming and outgoing liquidity in each channel
\end{enumerate}

\bigskip


\subsection*{Question 53}
Which service allows moving liquidity to the opposite side of a channel while retrieving funds on-chain?

\begin{enumerate}[label=\Alph*., leftmargin=*]
\item CoinSwap
\item Keysend
\item LNURL
\item Loop Out
\end{enumerate}

\bigskip


\subsection*{Question 54}
Why must a router maintain good liquidity distribution in its channels?

\begin{enumerate}[label=\Alph*., leftmargin=*]
\item To avoid opening new channels
\item To reduce transaction fees
\item To maximize the number of payments it can route
\item To limit the use of the Bitcoin blockchain
\end{enumerate}

\bigskip


\subsection*{Question 55}
What type of liquidity must a seller on the Lightning Network primarily manage to accept payments?

\begin{enumerate}[label=\Alph*., leftmargin=*]
\item Outgoing liquidity
\item Incoming liquidity
\item A high on-chain balance
\item A high total channel capacity
\end{enumerate}

\bigskip


\newpage
\section*{Answer Key}

\begin{enumerate}
\item \textbf{Question 1:} B
\\ \textit{The Lightning Network aims to enhance Bitcoin's scalability by enabling fast and inexpensive transactions through off-chain payment channels.}

\item \textbf{Question 2:} A
\\ \textit{The capacity of the channel corresponds to the total funds of the two participants and remains fixed to determine the possible exchanges.}

\item \textbf{Question 3:} D
\\ \textit{After sending 40,000 satoshis, Alice's balance in the channel is 60,000 satoshis.}

\item \textbf{Question 4:} D
\\ \textit{Each participant is limited by the amount of satoshis they hold in the channel for their transactions, ensuring that the funds are always covered.}

\item \textbf{Question 5:} D
\\ \textit{The Lightning Network enables fast and low-cost transactions through off-chain payment channels, which differs from regular transactions on the blockchain.}

\item \textbf{Question 6:} C
\\ \textit{The private key is used to sign transactions to satisfy a script and unlock the bitcoins.}

\item \textbf{Question 7:} A
\\ \textit{A UTXO, or Unspent Transaction Output, is an unspent transaction output, representing a certain amount of bitcoins that the holder can use.}

\item \textbf{Question 8:} C
\\ \textit{Like a banknote, a UTXO must be spent entirely in a transaction, and any excess is returned to the sender as change.}

\item \textbf{Question 9:} A
\\ \textit{A 2/2 multisignature address requires the signatures of both parties involved for the funds to be spent, which is used for the security of Lightning channels.}

\item \textbf{Question 10:} A
\\ \textit{The locking script imposes conditions, such as signing with a private key, to unlock and use the funds in a transaction.}

\item \textbf{Question 11:} D
\\ \textit{The opening of a channel begins with an exchange of messages between the two parties, where they share the amounts they wish to commit and their public keys.}

\item \textbf{Question 12:} D
\\ \textit{Alice creates a withdrawal transaction to retrieve her funds if a problem occurs with Bob, thus ensuring the security of her deposit.}

\item \textbf{Question 13:} A
\\ \textit{The channel is considered open once the deposit transaction has been included in a Bitcoin block and has reached a certain number of confirmations.}

\item \textbf{Question 14:} B
\\ \textit{Without Bob's signature on the withdrawal transaction, Alice cannot ensure the recovery of her funds in case of an issue, posing a risk of loss.}

\item \textbf{Question 15:} D
\\ \textit{The last step to open a channel is to publish the deposit transaction on the blockchain, which officially opens the channel.}

\item \textbf{Question 16:} C
\\ \textit{A commitment transaction represents the current distribution of funds in the channel after each Lightning transaction, but is not published on the blockchain as long as the channel remains open.}

\item \textbf{Question 17:} A
\\ \textit{Commitment transactions are built and signed for each change in allocation, but they are only eventually published on the blockchain upon channel closure.}

\item \textbf{Question 18:} B
\\ \textit{When Bob sends an invoice to Alice, she creates and signs a new commitment transaction to adjust the distribution of funds, thereby transferring 30,000 satoshis in the channel.}

\item \textbf{Question 19:} D
\\ \textit{Each transaction within the channel updates the channel's state with a new commitment transaction, reflecting the current distribution of funds between the two participants.}

\item \textbf{Question 20:} D
\\ \textit{The signed commitment transactions ensure that each participant can recover their funds by publishing the latest version in case the channel closes.}

\item \textbf{Question 21:} C
\\ \textit{The main risk is that a party may attempt to cheat by publishing an old commitment transaction, which could allow them to reclaim funds they no longer have.}

\item \textbf{Question 22:} C
\\ \textit{The timelock imposes a delay during which the cheater cannot retrieve their funds, thus giving the other party the opportunity to respond.}

\item \textbf{Question 23:} C
\\ \textit{The revocation key allows the aggrieved party to recover all the funds in the event of an attempted cheat by the other party.}

\item \textbf{Question 24:} A
\\ \textit{Exchanging secrets for previous transactions allows each participant to punish the other in case of an attempt to cheat by publishing an old transaction.}

\item \textbf{Question 25:} C
\\ \textit{The timelock and the revocation key deter participants from cheating by ensuring that any attempt to cheat results in an immediate loss of funds for the guilty party.}

\item \textbf{Question 26:} C
\\ \textit{There are three methods: cooperative closure, where the parties agree; forced closure without a response from one party; and closure with cheating, where a party attempts to publish an old transaction.}

\item \textbf{Question 27:} A
\\ \textit{Cooperative closure is the fastest and most economical method because both parties agree to close the channel, adjusting the fees according to market conditions.}

\item \textbf{Question 28:} C
\\ \textit{Forced closure is less desirable because it includes a time lock, which slows down the process, and the fees may be inappropriate since they were set at the time of the channel's creation.}

\item \textbf{Question 29:} B
\\ \textit{In the event of cheating, Bob can use the revocation key to recover all the funds from the channel, thus punishing Alice for her attempt.}

\item \textbf{Question 30:} D
\\ \textit{Forced closure is used as a last resort if one of the parties, for example, due to a technical problem, no longer responds to the other's messages.}

\item \textbf{Question 31:} B
\\ \textit{The Lightning Network enables the routing of a payment through several intermediary nodes, without requiring a direct channel between the sender and the recipient.}

\item \textbf{Question 32:} D
\\ \textit{The maximum amount that an indirect payment can transfer depends on the lowest liquidity available in the channels used, as each channel must adjust the funds on each side.}

\item \textbf{Question 33:} B
\\ \textit{Routing fees are used to compensate intermediary nodes that relay payments by offering their liquidity to enable the transfer of funds across the network.}

\item \textbf{Question 34:} C
\\ \textit{Onion routing means that each intermediate node only knows the previous and the next node, thus preserving the confidentiality of the final recipient.}

\item \textbf{Question 35:} C
\\ \textit{The sending node must know the network topology to calculate the route to the recipient, enabling it to choose the most optimal path.}

\item \textbf{Question 36:} C
\\ \textit{HTLCs (Hashed Time-Locked Contracts) ensure that each intermediary node transfers the payment without retaining the funds, as it can only unlock them by completing the transaction up to the final recipient.}

\item \textbf{Question 37:} A
\\ \textit{To receive the funds in an HTLC, the recipient must provide a specific secret that allows unlocking the payment secured by the hash.}

\item \textbf{Question 38:} A
\\ \textit{If an HTLC expires without the secret being revealed, the payment is canceled, and the funds are returned to the sender, thus preventing any intermediary node from retaining the funds.}

\item \textbf{Question 39:} A
\\ \textit{The timelock imposes a timeframe to complete an HTLC, allowing for the recovery of funds if the transaction does not succeed, by canceling the payment after this period.}

\item \textbf{Question 40:} C
\\ \textit{During a forced closure, pending HTLCs are represented as outputs in the commitment transactions, thus ensuring their execution or return according to the defined conditions.}

\item \textbf{Question 41:} D
\\ \textit{The sending node is responsible for calculating the entire route to the recipient, as intermediary nodes only know their immediate neighbor.}

\item \textbf{Question 42:} D
\\ \textit{The precise distribution of funds in a channel is only visible to the two participating nodes, preventing other nodes from estimating the available liquidity for routing.}

\item \textbf{Question 43:} B
\\ \textit{The total capacity of the channels can indicate a higher probability of success, as it suggests a greater availability of liquidity for transactions.}

\item \textbf{Question 44:} C
\\ \textit{By indicating nearby channels with liquidity or revealing private channels, Bob helps Alice to choose routes with a better probability of success.}

\item \textbf{Question 45:} B
\\ \textit{When a route fails, the sending node tries other alternative routes until a functional path is found or all options are exhausted.}

\item \textbf{Question 46:} C
\\ \textit{LNURL-withdraw allows users to withdraw funds from a service by scanning a QR code, thus simplifying the withdrawal process.}

\item \textbf{Question 47:} B
\\ \textit{The prefix 'lnbc' signifies that the invoice is for a Lightning transaction on the Bitcoin network, not on another network.}

\item \textbf{Question 48:} B
\\ \textit{LNURL is a communication protocol designed to simplify interactions, such as withdrawals and payments, between nodes and Lightning applications.}

\item \textbf{Question 49:} B
\\ \textit{Keysend allows the sender to generate a secret embedded in the HTLCs, enabling the recipient to unlock the funds without having issued an invoice.}

\item \textbf{Question 50:} C
\\ \textit{A Lightning invoice is a payment request, transmitted by the receiving node to the sending node to initiate a payment on the Lightning Network.}

\item \textbf{Question 51:} B
\\ \textit{By making a fictitious payment, a seller moves liquidity from their side to the opposite side of the channel, thereby freeing up incoming capacity to receive payments.}

\item \textbf{Question 52:} D
\\ \textit{The triangle opening technique allows each node to achieve balanced liquidity, both incoming and outgoing, by forming a loop of channels among several participants.}

\item \textbf{Question 53:} D
\\ \textit{The Loop Out service from Lightning Labs enables the transfer of funds from a Lightning channel to the on-chain blockchain, thereby optimizing receiving capacity while retrieving bitcoins on the main chain.}

\item \textbf{Question 54:} C
\\ \textit{A router aims to route as many payments as possible to collect fees, which requires a good distribution of liquidity on both sides of the channels to ensure routing capacity.}

\item \textbf{Question 55:} B
\\ \textit{To receive payments, a seller needs to have incoming liquidity, meaning that funds must be available on the opposite side of the channel for them to accept payments.}


\end{enumerate}

\end{document}
